% 原书分章 PDF 第 10—11 页：第 23.6 节。
\section{讨论}

我们研究了最优资源约束价值的估计问题。在因果假设下，这一参数可以识别为：在依赖已测量协变量的个体化处理规则下，在总体中最多只有 $\kappa$ 的比例能够接受处理这一约束下，所能达到的最大总体平均结局。我们还给出了资源约束下一个最优随机规则的显式表达式。

此外，在“某个具有正概率的协变量分层中的处理效应不恰好等于 $\tau_0$”这一关键假设下，我们推导了最优 R-C 价值的典范梯度。典范梯度在构造渐近线性估计量中起关键作用。我们发现，当资源约束起作用时，即 $\tau_0>0$ 时，与最优无约束价值的典范梯度相比，最优 R-C 价值的典范梯度多出一个分量。

\Needspace{10\baselineskip}
\begin{sourcebox}
我们提出了最优 R-C 价值的 TMLE。这个估计量的设计使其求解典范梯度某个估计的经验均值方程。由此很快得出了使估计量为 RAL、且在所有这类 RAL 估计量中有效的条件。所有这些结果都依赖于以下条件：对于具有正概率的协变量分层，处理效应不恰好等于 $\tau_0$。当协变量连续且约束起作用时，这个假设比通常的非例外分布假设更合理，因为处理效应集中在任意的（由 $\kappa$ 决定的）$\tau_0>0$ 上可能并不常见。注意，这个拟非例外分布假设（pseudo-nonexceptional law assumption）已经蕴含，最优随机规则几乎处处等于最优确定性规则。
\end{sourcebox}

实践中遇到的某些资源约束，未必具有 $\E[d(U,W)]$ 小于或等于 $\kappa$ 的形式。例如，向个体分配处理的成本可能随协变量的取值而变化。如果 $c:\mathcal W\to[0,\infty)$ 为成本函数，那么约束可以具有 $\E[c(W)d(U,W)]\leq\kappa$ 的形式。对于满足 $\bar Q_{b,0}(w)\ne\tau_0c(w)$ 或 $c(w)=0$ 的 $w$，最优规则具有 $(u,w)\mapsto I(\bar Q_{b,0}(w)>\tau c(w))$ 的形式；剩余资源则在其余所有 $w$ 之间均匀随机分配。我们把对这一更一般资源约束问题的进一步研究留待未来工作。

我们没有研究向总体分配有限资源所涉及的伦理问题。关于如何向总体分配有限处理资源才合适，争论仍在继续（例如，HIV/AIDS 治疗中的讨论见 Brock 与 Wikler，2009；Macklin 与 Cowan，2012；Singh，2013）。显然，任何研究者都需要考虑某些资源分配方案所涉及的伦理问题。我们的方法在一种特定的功利主义意义下是最优的（针对感兴趣的结局，使总体平均结局的期望最大），并得出一种处理策略：让那些在感兴趣结局上预期从处理中获益最多的人接受处理。必须谨慎，确保感兴趣的结局确实反映了最重要的公共卫生影响。与无约束的个体化医学不同，不恰当地向某个分层开出处方，也会影响该分层之外的个体，即那些因处理资源不可获得而未接受处理的人。我们把进一步的伦理探讨留给相关专家。未来值得研究的问题是，是否存在某些情形，可以通过变换结局或为优化问题增加约束，使本章考虑的统计问题遵循这些情形中的伦理准则。

我们试图把以往估计最优个体化处理策略价值的工作，推广到处理资源有限的情形，即不可能对总体中的所有人实施处理的情形。本章结果应当能使最优个体化处理策略应用于许多新的重要问题。本章经许可包含了 Luedtke 与 van der Laan（2016c）的部分内容。
